\section{Discussion}

The case study shows why HIL-RL trigger evaluation should be stricter than a single learned-trigger-versus-random comparison. The original LV-VoI result looked plausible until the random baseline was cost matched. After the random budget sweep, random b350 was both more successful and less expensive than LV-VoI scale3. This reversal is the central empirical lesson of the paper: intervention-trigger quality cannot be judged from nominal budget settings alone, because realized human-step cost determines the relevant comparison. A trigger that appears selective at one baseline budget may be dominated once the random family is placed on the same success-cost axis.

Repeated checkpoint evaluation is the second part of the same discipline. In noisy off-policy robot learning, a final checkpoint or a single restored-best evaluation can make a policy look better or worse than its expected behavior. The current evidence package therefore reports repeated checkpoint success counts and binomial intervals rather than single point estimates. This does not remove all variance, but it makes the uncertainty visible enough to prevent a small or lucky evaluation from carrying a paper-level claim.

The trace analysis is equally important because it changes what kind of repair is justified. The failed trigger was not simply asking far from the object: LV-VoI starts were closer to the cube than random starts, yet they were more frequent and less cost effective. The score-floor repair then showed that a mechanically valid threshold can still fail scientifically if it does not reduce accepted starts enough to recover the frontier. These findings argue against quick heuristic repairs based only on one intuitive failure story. They support a diagnostic workflow in which mechanism hypotheses are checked against start-level timing, geometry, and score distributions before a new trigger is scaled.

The current evidence therefore supports a protocol contribution rather than an algorithmic superiority claim. This distinction is important for a journal-style paper. A dominated trigger result can still be useful when it exposes a systematic evaluation failure mode and turns that failure into a reusable reporting standard. FORESIGHT-HIL should be presented as the audited case study that motivated the protocol, not as a method that currently beats random intervention. The paper's positive contribution is the evaluation discipline: cost-matched random families, repeated checkpoint estimates, trace-level intervention diagnostics, and same-seed stop rules for lightweight repairs.

Several limitations remain. The intervention source is a scripted privileged-state oracle, not a human teleoperator, so the results do not measure real human workload, delay, fatigue, or interface effects. The strongest evidence is concentrated on robosuite Lift, while Stack currently functions as boundary evidence showing that online intervention can underperform a strong no-online baseline. The paper also does not yet provide a redesigned trigger that recovers the success-cost frontier. These limitations should remain explicit because they protect the central claim: before HIL-RL papers claim trigger superiority, they should show that the claim survives cost matching, repeated evaluation, and trace-level budget analysis.

The practical implication is a stop-and-redesign rule. If a trigger variant is dominated by same-seed random on both repeated success and realized human cost, it should not be expanded only because it is intuitively appealing. Future method work should begin from the trace findings rather than from the earlier unsupported superiority story. A stronger trigger may need explicit budget control, calibrated trigger scores, task-phase awareness, or contact-aware intervention logic, but those ideas should enter as new experiments only after the current diagnostic paper spine is stable.
